\subsubsection{经由 Fourier 变换的构造}\label{sss:Four} 由 $\pi$ 的半小性，复形 $\calS=\bR\pi_{*}\Ql[\dim\sN]$ 也是反常层。然而它不是从 $\calN$ 的某个开子集作中间扩充得来的。仿射空间上 $\Gm$ 等变层有 Fourier 变换的概念 \cite{Laumon}。可以证明 $\calS$ 同构于 $\calS_{\frg}$ 的 Fourier 变换，反之亦然。于是由 Fourier 变换的函子性，$\calS_{\frg}$ 上的 $W$ 作用诱导出 $\calS$ 上的 $W$ 作用。取 $\calS$ 在 $e$ 处的茎，就得到 $W$ 在 $\cohog{*}{\sB_{e}}$ 上的作用。按 \S\ref{sss:two actions} 的约定，这个作用是规范化的。

\subsubsection{经由对应的构造} 考虑 Steinberg 簇 $\St_{\frg}=\tg\times_{\frg}\tg$，它分类三元组 $(X,B_{1},B_{2})$，其中 $B_{1},B_{2}$ 是 $G$ 的 Borel 子群，且 $X\in\Lie B_{1}\cap\Lie B_{2}$。$\St_{\frg}$ 的不可约分支以 Weyl 群的元素为指标：对 $w\in W$，令 $\St_{w}$ 为 $w$ 在 $\tg^{\rs}$ 上作用的图像之闭包，则 $\St_{w}$ 是 $\St$ 的不可约分支，且这些穷尽 $\St$ 的全部不可约分支。上同调对应的形式体系使我们能从每个 $\St_{w}$ 得到复形 $\calS_{\frg}=\bR\pi_{\frg,*}\Ql[\dim\frg]$ 的一个自同态。要证明这些自同态合起来构成 $W$ 在 $\calS_{\frg}$ 上的作用，并不显然。论证的关键仍是映射 $\pi_{\frg}$ 的小性。定义了 $\calS_{\frg}$ 上的 $W$ 作用之后，可以用 \S\ref{sss:mid} 的办法，把 $W$ 在茎 $\calS_{\frg, e}$ 上的作用与符号表示作张量，也可以用 \S\ref{sss:Four} 的 Fourier 变换，来定义 $\cohog{*}{\sB_{e}}$ 上的 Springer 作用。关于这一构造的讨论参见 \cite[Remark 3.3.4]{GS}；亦可参见 \cite[\S3.4]{CG}，那里用 $\St_{w}$ 在幂零 Steinberg 簇 $\tN\times_{\sN}\tN$ 中的极限给出了类似但不同的构造。

注意上述三种构造都能证明 $\End(\calS)\cong\Ql[W]$，从而给出定理 \ref{th:BM} 的证明。

\subsubsection{经由单值的构造} 我们概述 Slodowy \cite[\S4]{Slo4} 的一个构造，它适用于 $k=\CC$。当时猜想这一构造给出的 $W$ 在 $\cohog{*}{\sB_{e}}$ 上的作用与定理 \ref{th:Spr corr} 中的相同。Rossmann 在 \cite[\S2]{Rossmann} 中给出了类似的构造，他把自己的作用等同于 Kazhdan 与 Lusztig 在 \cite{KL80} 中构造的作用，而后者已知与 Springer 的作用相同。因此这些构造给出的 $W$ 作用都与定理 \ref{th:Spr corr} 中的相同。

设 $e\in\calN$，设 $S_{e}\subset \frg$ 是 $e$ 的轨道的横截片。选取包含 $e$ 的 $\sl_{2}$ 三元组 $(e,h,f)$ 后，有这样的横截片的标准取法：$S_{e}=e+\frg_{f}$，其中 $\frg_{f}$ 是 $f$ 在 $\frg$ 中的中心化子。现考虑下图，其中除最右方块外各方块都是笛卡儿的
\begin{equation}\label{Slo}
\xymatrix{\sB_{e}\ar[d]\ar@{^{(}->}[r] & \tilS^{\nil}_{e}\ar[d]\ar@{^{(}->}[r] & \tilS_{e}\ar[d]\ar@{^{(}->}[r] & \tg\ar[d]^{\pi_{\frg}}\ar[r]^{\wt\chi} & \frt\ar[d]\\
\{e\}\ar@{^{(}->}[r] & S^{\nil}_{e}\ar@{^{(}->}[r] & S_{e}\ar@{^{(}->}[r] & \frg\ar[r]^{\chi} & \frc}
\end{equation}
这里最右方块是 \eqref{tgg}，$S^{\nil}_{e}=S_{e}\cap\sN$，而 $\tilS_{e}$ 与 $\tilS^{\nil}_{e}$ 分别是 $S_{e}$ 与 $S^{\nil}_{e}$ 在 $\pi_{\frg}$ 下的原像。取 $e$ 处的小球 $V_{e}\subset S_{e}$，再取 $0\in\frc$ 处更小的球 $V_{0}\subset \frc$。令 $U_{e}=V_{e}\cap \chi^{-1}(V_{0})\subset S_{e}$。则图 \eqref{Slo} 限制为图
\begin{equation}\label{Slo U}
\xymatrix{\sB_{e}\ar[d]\ar@{^{(}->}[r] & \tilU^{\nil}_{e}\ar[d]\ar@{^{(}->}[r] & \tilU_{e}\ar[d]\ar[r]^{\wt\chi_{e}} & \tilV_{0}\ar[d]\\
\{e\}\ar@{^{(}->}[r] & U^{\nil}_{e}\ar@{^{(}->}[r] & U_{e}\ar[r] & V_{0}}
\end{equation}
这里 $U^{\nil}_{e}=U_{e}\cap\sN$，而 $\tilV_{0}, \tilU_{e}$ 与 $\tilU^{\nil}_{e}$ 分别是 $V_{0}, U_{e}$ 与 $U^{\nil}_{e}$ 在 \eqref{Slo} 中竖直映射下的原像。关键的拓扑事实有两条
\begin{itemize}
\item 包含 $\sB_{e}\incl \tilU^{\nil}_{e}$ 有形变收缩，因此是同伦等价；
\item 映射 $\wt\chi_{e}: \tilU_{e}\to \tilV_{0}$ 是可平凡化的纤维丛（微分拓扑意义下）。
\end{itemize}
现在，由第二条事实，$\wt\chi_{e}$ 的一般纤维带有一个同伦意义下的 $W$ 作用：因为 \eqref{Slo U} 的最右方块在 $V_{0}\cap\frc^{\rs}$ 上方是笛卡儿的，且 $\tilV_{0}\to V_{0}$ 在 $V_{0}\cap\frc^{\rs}$ 上方是 $W$-挠子。由第一条事实，$\sB_{e}$ 与 $\tilU^{\nil}_{e}=\wt\chi^{-1}_{e}(0)$ 同伦等价；又因 $\wt\chi_{e}$ 是纤维丛，$\sB_{e}$ 也与 $\wt\chi_{e}$ 的一般纤维同伦等价。把这些事实合起来，就得到 $W$ 在 {\em $\sB_{e}$ 的同伦型}上的作用，这比 $W$ 作用在 $\sB_{e}$ 的上同调上更强。此构造的一个推论是：定理 \ref{th:Spr corr} 中 $\cohog{*}{\sB_{e}}$ 上的 $W$ 作用保持环结构。

\subsubsection{在定理 \ref{th:BM} 之下证明定理 \ref{th:Spr corr}}
把反常层 $\calS$ 按 $W$ 作用分解为同型分量
\begin{equation*}
\calS=\bigoplus_{\chi\in\Irr(W)}V_{\chi}\otimes\calS_{\chi}
\end{equation*}
其中 $V_{\chi}$ 是 $W$ 通过不可约表示 $\chi$ 作用的那个空间，而 $\calS_{\chi}=\Hom_{W}(V_{\chi},\calS)$ 是 $\sN$ 上的反常层。由 $\End(\calS)\cong\Ql[W]$ 可得每个 $\calS_{\chi}$ 都非零，且
\begin{equation}\label{HomSS}
\Hom(\calS_{\chi},\calS_{\chi'})=\begin{cases}\Ql & \chi=\chi' \\ 0 & \text{其他情形}\end{cases}
\end{equation}
分解定理 \cite[Th 6.2.5]{BBD} 蕴含每个 $\calS_{\chi}$ 都是半单反常层。因此 \eqref{HomSS} 说明 $\calS_{\chi}$ 是单的反常层。于是 $\calS_{\chi}$ 形如 $\IC(\ov\calO, \calL)$，其中 $\calO\subset\sN$ 是幂零轨道，$\calL$ 是 $\calO$ 上不可约 $G$ 等变局部系统。再者，由于 $\chi\neq\chi'$ 时 $\Hom(\calS_{\chi},\calS_{\chi'})=0$，这些单反常层 $\{\calS_{\chi}\}_{\chi\in\Irr(W)}$ 互不同构。把 \eqref{Spr corr} 的右边理解为幂零轨道上不可约 $G$ 等变局部系统的同构类，就证明了定理 \ref{th:Spr corr} 的第 (3) 部分。若 $\calS_{\chi}=\IC(\ov\calO,\calL)$ 且 $e\in\calO$，则 $\pi$ 的半小性使我们可以把茎 $\calL_{e}$ 与 $\cohog{2d_{e}}{\sB_{e}}$ 的一个 $A_{e}$ 同型子空间等同起来。这就证明了定理 \ref{th:Spr corr} 的第 (2) 部分。\qed

下面再举几个 Springer 对应的例子。

\subsubsection{$A$ 型}\label{sss:A Spr corr} 当 $G=\SL_{n}$ 时，所有 $A_{e}$ 都平凡。Springer 对应给出 $W=S_{n}$ 的不可约表示与 $\frg=\sl_{n}$ 的幂零轨道之间的双射，两者都由 $n$ 的分拆参数化。在 \S\ref{sss:Spr type A} 中已经看到：若 $e$ 的 Jordan 型为 $\l$，则最高维上同调 $\cohog{2d_{e}}{\sB_{e}}$ 有一组以 $\l$ 的标准 Young 表为指标的基，而后者也给出 $S_{n}$ 对应于分拆 $\l$ 的不可约表示的一组基。例如，对 $G=\SL_{3}$，Springer 对应为
\begin{itemize}
\item 平凡表示 $\bij$ 正则轨道，分拆 $3=3$；
\item 二维表示 $\bij$ 次正则轨道，分拆 $3=2+1$；
\item 符号表示 $\bij$ $\{0\}$，分拆 $3=1+1+1$。
\end{itemize}

\subsubsection{次正则轨道与反射表示}\label{sss:subreg Spr corr} 考虑 $e$ 是次正则幂零元的情形。此时分支群 $A_{e}$ 可以与 \S\ref{sss:subreg} 中引入的 $G'$ 的 Dynkin 图的自同构群等同（见 \cite[\S7.5, Proposition]{Slodowy}）。把 $\sB_{e}$ 的不可约分支与 $G'$ 的 Dynkin 图的顶点等同之后，$A_{e}$ 在 $\cohog{2}{\sB_{e}}$ 上的作用以「重排由不可约分支给出的基」的方式实现，与它在 $G'$ 的 Dynkin 图上的作用相同。例如当 $G=G_{2}$ 时，可写 $\sB_{e}=C_{1}\cup C_{2}\cup C_{3}\cup C_{4}$，其中 $C_{1}$、$C_{2}$、$C_{3}$ 各自与 $C_{4}$ 交于一点，其余两两不交。群 $A_{e}$ 同构于 $S_{3}$，它在 $\cohog{2}{\sB_{e}}$ 上的作用固定 $C_{4}$ 的基本类，并置换 $C_{1}, C_{2}, C_{3}$ 的基本类。

注意 $\cohog{2}{\sB_{e}}^{A_{e}}$ 的维数总是 $r$，即 $G$ 的秩。事实上，作为 $W$ 模，$\cohog{2}{\sB_{e}}^{A_{e}}$ 同构于 $W$ 在 $\frt^{*}$ 上的反射表示。换言之，在 Springer 对应之下，数对 $(e=\text{次正则}, \rho=1)$ 对应 $W$ 的反射表示。
